\section{Homomorphism of complexes}\label{sec homs}

\subsection[Lie algebra $\mathfrak{sl}_2(U)$]{Lie algebra $\boldsymbol{\mathfrak{sl}_2(U)}$}\label{conf block}
Recall that $\{z_1,\dots,z_n,z_{n+1}=\infty\}$ are pairwise distinct points of the complex projective line~${\mathbb P}^1$ and $U={\mathbb P}^1 -\{z_1,\dots,z_n,z_{n+1}\}$. Fix local coordinates $t-z_1,\dots, t-z_n,1/t$ on ${\mathbb P}^1$ at these points, respectively.
Let $\mathfrak{sl}_2(U)$ be the Lie algebra of $\mathfrak{sl}_2$-valued rational functions on ${\mathbb P}^1$ regular on $U$, with the pointwise bracket. Thus, an element of $\mathfrak{sl}_2(U)$ has the form $e\otimes u_1 + h\otimes u_2+ f\otimes u_3$ with $u_i\in\Omega^0(U)$, and the bracket is defined by the formula $[x\otimes u_1, y\otimes u_2] = [x,y]\otimes (u_1u_2)$.

\subsection[$\mathfrak{sl}_2(U)$-modules]{$\boldsymbol{\mathfrak{sl}_2(U)}$-modules}
We say that an $\widehat{{\mathfrak{sl}_2}}$-module $W$ has the finiteness property, if for any $w\in W$ and $ x\in\mathfrak{sl}_2$, we have $xT^j \cdot w=0$ for all $j \gg 0$. For example, the contragradient Verma module has the finiteness property.

Let $W_1,\dots,W_{n+1}$ be $\widehat{{\mathfrak{sl}_2}}$-modules with the finiteness property. Then the Lie algebra~$\mathfrak{sl}_2(U)$ acts on $W_1\otimes\dots\otimes W_{n+1}$ by the formula
\begin{gather*}
x\otimes u \cdot (w_1\otimes\dots\otimes w_{n+1}) = \big([x\otimes u(t)]^{(z_1)} w_1\big)
\otimes w_2\otimes\dots\otimes w_{n+1}+ \cdots \nonumber\\
\hphantom{x\otimes u \cdot (w_1\otimes\dots\otimes w_{n+1}) =}{} + w_1\otimes\dots\otimes w_{n-1}\otimes \big([x\otimes u(t)]^{(z_n)} w_n\big)\otimes w_{n+1} \nonumber\\
\hphantom{x\otimes u \cdot (w_1\otimes\dots\otimes w_{n+1}) =}{} + w_1\otimes\dots\otimes w_n\otimes \big(\pi([x\otimes u(t)]^{(\infty)}) w_{n+1}\big), %\label{action}
\end{gather*}
where for $x\otimes u \in\mathfrak{sl}_2(U)$ the symbol $[x\otimes u(t)]^{(z_j)}$ denotes the Laurent expansion of $x\otimes u$ at $t=z_j$ and $[x\otimes u(t)]^{(\infty)}$ denotes the Laurent expansion at $t=\infty$; the symbol $\pi$ in the last term denotes the $\widehat{{\mathfrak{sl}_2}}$-automorphism defined in Section \ref{aut}.

The finiteness property of the tensor factors ensures that the actions of the Laurent series are well-defined.

The $\widehat{{\mathfrak{sl}_2}}$-action gives us a map
\begin{gather}\label{mult}
\mu\colon \ \mathfrak{sl}_2(U)\otimes \big(\otimes_{j=1}^{n+1} W_j\big) \to \otimes_{j=1}^{n+1} W_j .
\end{gather}

\subsection{Chain complex}
For a Lie algebra ${\mathfrak g}$ and a ${\mathfrak g}$-module $W$ we denote by $C_{\bullet}({\mathfrak g}, W)$ the standard chain complex of ${\mathfrak g}$ with coefficients in $W$, where
\begin{gather*}
 C_p({\mathfrak g}, W) = \wedge ^p{\mathfrak g} \otimes W, \\
{\rm d}(g_p\wedge \dots\wedge g_1\otimes w) = \sum_{i=1}^p(-1)^{i-1} g_p\wedge \dots\wedge \widehat{g_i}\wedge\dots \wedge g_1\otimes g_iw \\
\hphantom{{\rm d}(g_p\wedge \dots\wedge g_1\otimes w) =}{} + \sum_{1\leq i<j\leq p}(-1)^{i+j} g_p\wedge \dots\wedge \widehat{g_j}\wedge\dots\wedge \widehat{g_i}\wedge\dots \wedge g_1\otimes [g_j,g_i]w.
\end{gather*}

\subsection{Two complexes}

\subsubsection{} Let $m_1,\dots,m_n, k\in{\mathbb C} $, $k+2\ne 0$. Define $m_{n+1}=m_1+\dots+m_n-2$. For $j=1,\dots, n+1$, let $V_j$ be the $\widehat{{\mathfrak{sl}_2}}$ Verma module $V(m_j,k-m_j)$ and $V_j^*$ the corresponding contragradient Verma module. Consider the chain complex $C_{\bullet}\big(\mathfrak{sl}_2(U), \otimes_{j=1}^{n+1}V_j^*\big)$ and its last two terms{\samepage
\begin{gather*}%\label{l2t}
\to\ \mathfrak{sl}_2(U) \otimes \big({\otimes}_{j=1}^{n+1}V_j^*)\overset{{\rm d}}{\longrightarrow} \otimes_{j=1}^{n+1}V_j^* \to 0,
\end{gather*}
where ${\rm d}=\mu$, see formula (\ref{mult}).}

We assign degree $0$ to the term $\otimes_{j=1}^{n+1}V_j^*$ of this complex and assign degree $1$ to the differential~${\rm d}$, so that the whole complex sits in the non-positive area.

\subsubsection{} Consider the twisted de Rham complex in (\ref{dr com}) corresponding to $\kappa=k+2$ with degrees shifted by~1, namely, the complex $\Omega^{\bullet}(U)[1]$,
\begin{gather*}%\label{cder}
 0\to\Omega^0(U) \overset{\partial}{\longrightarrow} \Omega^1(U)\to 0,
\end{gather*}
where the shift $[1]$ means that we assign degree $p-1$ to the term $\Omega^p(U)$.

\subsection{Construction}\label{sec constr}
Define a linear map
\begin{gather*}
\eta^1\colon \ \Omega^1(U)\longrightarrow \otimes_{j=1}^{n+1} V_j^*
\end{gather*}
by the formulas
\begin{gather}
\frac{{\rm d}t}{(t-z_m)^{a}} \mapsto -\kappa (v_1)^*\otimes\dots\otimes \left(\frac{f}{T^{a-1}}v_m\right)^* \otimes\dots\otimes (v_{n+1})^*, \label{eta 1f} \\
t^{a-1}{\rm d}t \mapsto \kappa (v_1)^*\otimes\dots\otimes (v_n)^* \otimes \left(\frac{e}{T^{a}}v_{n+1}\right)^*, \label{eta 1i}
\end{gather}
for $a> 0$. Define a linear map
\begin{gather*}
\eta^0\colon \ \Omega^0(U)\longrightarrow \mathfrak{sl}_2(U)\otimes \big({\otimes}_{j=1}^{n+1} V_j^*\big)
\end{gather*}
by the formulas
\begin{gather}
\frac{1}{(t-z_m)^a} \mapsto \frac{f}{(t-z_m)^a}\otimes (v_1)^*\otimes\dots\otimes (v_{n+1})^*\nonumber\\
\hphantom{\frac{1}{(t-z_m)^a} \mapsto}{} -\sum_{l=1}^a\bigg[\frac{e}{(t-z_m)^l}\otimes (v_1)^*\otimes \dots\otimes 2\sum_{i+j=a-l \atop i\geq j\geq 0} \left(\frac{f}{T^i}\frac{f}{T^j}v_m\right)^*\otimes\dots\otimes (v_{n+1})^* \nonumber \\
\hphantom{\frac{1}{(t-z_m)^a} \mapsto}{} {} +\frac{h}{(t-z_m)^l}\otimes (v_1)^*\otimes \dots\otimes \left(\frac{f}{T^{a-l}}v_m\right)^*\otimes \dots\otimes (v_{n+1})^* \bigg],\label{eta 0f}
\end{gather}
for $a>0$;
\begin{gather}
t^a \mapsto ft^a\otimes (v_1)^*\otimes\dots\otimes (v_{n+1})^*\nonumber \\
\hphantom{t^a \mapsto}{} - \sum_{l=0}^{a-2}\bigg[et^l\otimes (v_1)^*\otimes\dots\otimes
(v_n)^*\otimes 2\sum_{i+j=a-l,\atop i\geq j\geq 1} \left(\frac{e}{T^i} \frac{e}{T^j}v_{n+1}\right)^* \nonumber \\
\hphantom{t^a \mapsto}{} + ht^{l+1}\otimes (v_1)^*\otimes\dots\otimes (v_n)^*\otimes \left(\frac{e}{T^{a-l-1}}v_{n+1}\right)^* \bigg],\label{eta 0i}
\end{gather}
for $a\geq 0$.

\begin{thm}[{\cite[Theorem 5.12]{SV2}}]\label{thm hom} Formulas \eqref{eta 1f}--\eqref{eta 0i} define a homomorphism of complexes $\eta\colon \Omega^{\bullet}(U)[1] \to C_{\bullet}\big(\mathfrak{sl}_2(U);\otimes_{j=1}^{n+1} V_j^*\big)$, namely we have
\begin{gather*}%\label{homom}
{\rm d} \eta^0 = \eta^1\partial.
\end{gather*}
The homomorphism is injective.
\end{thm}

Theorem \ref{thm hom} was announced in \cite{SV2}. Here is a proof of the theorem.

\begin{proof} First we calculate $ \eta^1(\partial((t-z_p)^{-a}))$,
\begin{gather*}
\frac 1{(t-z_p)^{a}} \overset{\partial} \mapsto -\frac{1}{\kappa}(m_p+a\kappa) \frac{{\rm d}t}{(t-z_p)^{a+1}} +\frac{1}{\kappa} \sum_{k=1}^a \sum_{j\neq p}\frac{m_j}{(z_j-z_p)^k}
\frac{{\rm d}t}{(t-z_p)^{a+1-k}} \nonumber \\
\hphantom{\frac 1{(t-z_p)^{a}} \overset{\partial} \mapsto}{} - \frac{1}{\kappa} \sum_{j\neq p} \frac{m_j}{(z_j-z_p)^a}\frac{{\rm d}t}{t-z_j} \nonumber \\
\hphantom{\frac 1{(t-z_p)^{a}}}{}\overset{\eta^1} \mapsto (m_p+\kappa a) (v_1)^* \otimes\dots\otimes\left(\frac{f}{T^a}v_p \right)^* \otimes\dots\otimes (v_{n+1})^* \nonumber \\
\hphantom{\frac 1{(t-z_p)^{a}} \overset{\partial} \mapsto}{} - \sum_{k=1}^a \sum_{j \neq p}\frac{ m_j}{(z_j-z_p)^k} (v_1)^* \otimes\dots\otimes \left(\frac{f}{T^{a-k}} v_p \right)^* \otimes\dots\otimes (v_{n+1})^* \nonumber\\
\hphantom{\frac 1{(t-z_p)^{a}} \overset{\partial} \mapsto}{} + \sum_{j\neq p}\frac{ m_j}{(z_j-z_p)^a} (v_1)^* \otimes\dots\otimes (fv_j)^* \otimes\dots\otimes (v_{n+1})^*. \nonumber
\end{gather*}
Then we calculate ${\rm d}\big(\eta^0((t-z_p)^{-a})\big)$,
\begin{gather*}
\frac1{(t-z_p)^{a}} \overset{\eta^0}\mapsto \frac{f}{(t-z_p)^a} \otimes (v_1)^* \otimes\dots\otimes (v_{n+1})^* \nonumber \\
\hphantom{\frac1{(t-z_p)^{a}} \overset{\eta^0}\mapsto}{} - \sum_{l=1}^a \bigg[ \frac{h}{(t-z_p)^l} \otimes (v_1)^* \otimes\dots\otimes \left( \frac{f}{T^{a-l}} v_p \right)^* \otimes\dots\otimes (v_{n+1})^*
\nonumber \\
\hphantom{\frac1{(t-z_p)^{a}} \overset{\eta^0}\mapsto}{} + \frac{e}{(t-z_p)^l} \otimes 2 \sum_{i+j = a-l \atop i \geq j \geq 0} (v_1)^* \otimes\dots\otimes \left(\frac{f}{T^i} \frac{f}{T^j} v_p \right)^* \otimes\dots\otimes (v_{n+1})^* \bigg] \nonumber \\
\hphantom{\frac1{(t-z_p)^{a}}}{}
\overset{\mu} \mapsto (v_1)^* \otimes\dots\otimes \bigg[ (m_p + a(k+2)) \left(\frac{f}{T^a} v_p \right)^* \nonumber \\
\hphantom{\frac1{(t-z_p)^{a}} \overset{\eta^0}\mapsto}{} + \sum_{l=1}^a \bigg[\frac{h}{T^l}\left(\frac{f}{T^{a-l}}v_p\right)^*+
2\frac{e}{T^l}\sum_{i+j=a-l \atop i\geq j\geq 0 } \left(\frac{f}{T^i}\frac{f}{T^j} v_p\right)^*\bigg] \bigg] \otimes\dots\otimes (v_{n+1})^* \nonumber \\
\hphantom{\frac1{(t-z_p)^{a}} \overset{\eta^0}\mapsto}{} + \sum_{j\neq p} \frac{m_j}{(z_j-z_p)^a} (v_1)^* \otimes\dots\otimes (fv_j)^* \otimes\dots\otimes (v_{n+1})^* \nonumber \\
\hphantom{\frac1{(t-z_p)^{a}} \overset{\eta^0}\mapsto}{}
 - \sum_{l=1}^a \bigg[ (v_1)^* \otimes\dots\otimes \frac{h}{T^l}\left(\frac{f}{T^{a-l}}v_p \right)^* \otimes\dots\otimes (v_{n+1})^* \nonumber \\
\hphantom{\frac1{(t-z_p)^{a}} \overset{\eta^0}\mapsto}{} + (v_1)^* \otimes\dots\otimes \sum_{i+j=a-l \atop i\geq j\geq 0} 2\frac{e}{T^l} \left(\frac{f}{T^i}\frac{f}{T^j}v_p\right)^*\otimes\dots\otimes (v_{n+1})^*
\nonumber \\
\hphantom{\frac1{(t-z_p)^{a}} \overset{\eta^0}\mapsto}{} + \sum_{j \neq p} \frac{m_j}{(z_j-z_p)^l} (v_1)^* \otimes\dots\otimes \left(\frac{f}{T^{a-l}} v_p \right)^* \otimes\dots\otimes (v_{n+1})^* \bigg] \nonumber \\
\hphantom{\frac1{(t-z_p)^{a}}}{} = (\kappa a + m_p) (v_1)^* \otimes\dots\otimes\left(\frac{f}{T^a}v_p \right)^* \otimes\dots\otimes (v_{n+1})^* \nonumber \\
\hphantom{\frac1{(t-z_p)^{a}} \overset{\eta^0}\mapsto}{} - \sum_{l=1}^a \sum_{j \neq p} \frac{m_j}{(z_j-z_p)^l} (v_1)^* \otimes\dots\otimes \left(\frac{f}{T^{a-l}} v_p \right)^* \otimes\dots\otimes
(v_{n+1)}^* \nonumber \\
\hphantom{\frac1{(t-z_p)^{a}} \overset{\eta^0}\mapsto}{} + \sum_{j\neq p} \frac{ m_j}{(z_j-z_p)^a} (v_1)^* \otimes\dots\otimes (fv_j)^* \otimes\dots\otimes (v_{n+1})^*.
\end{gather*}
In this calculation we use formula (\ref{id A}) to express the action of $\frac{f}{T^a}$ on $(v_{p})^*$. These formulas show that ${\rm d}\big(\eta^0((t-z_p)^{-a})\big) = \eta^1\big(\partial((t-z_p)^{-a})\big)$.

Now we calculate $\eta^1(\partial(t^a))$,
\begin{gather*}
t^a \overset{\partial}{\mapsto} \frac{1}{\kappa}\bigg(a\kappa - \sum_{j=1}^n m_j\bigg) t^{a-1} {\rm d}t - \frac{1}{\kappa} \sum_{s=1}^{a-1} \sum_{j=1}^n m_j z_j^s
t^{a-s-1} {\rm d}t - \frac{1}{\kappa} \sum_{j=1}^n m_j z_j^a \frac{{\rm d}t}{t-z_j} \nonumber \\
\hphantom{t^a}{} \overset{\eta^1}{\mapsto} \bigg(a\kappa - \sum_{j=1}^n m_j\bigg) (v_1)^* \otimes \dots \otimes (v_n)^* \otimes \left(\frac{e}{T^a} v_{n+1}\right)^*\nonumber \\
\hphantom{t^a \overset{\partial}{\mapsto}}{}- \sum_{s=1}^{a-1} \sum_{j=1}^n m_j z_j^s (v_1)^* \otimes \dots \otimes (v_n)^* \otimes \left(\frac{e}{T^{a-s}} v_{n+1}\right)^* \nonumber \\
\hphantom{t^a \overset{\partial}{\mapsto}}{} + \sum_{j=1}^n m_j z_j^a (v_1)^* \otimes \dots \otimes (f v_j)^* \otimes\dots\otimes (v_{n+1})^*.
\end{gather*}

Then we calculate ${\rm d}\big(\eta^0(t^a)\big)$,
\begin{gather*}
t^a \overset{\eta^0}\mapsto ft^a\otimes (v_1)^*\otimes\dots\otimes (v_{n+1})^* - \sum_{l=0}^{a-2}\bigg[ht^{l+1}\otimes (v_1)^*\otimes\dots\otimes (v_n)^*\otimes \left(\frac{e}{T^{a-l-1}}v_{n+1}\right)^* \nonumber \\
\hphantom{t^a \overset{\eta^0}\mapsto}{} + et^l\otimes (v_1)^*\otimes\dots\otimes (v_n)^*\otimes 2\sum_{i+j=a-l\atop i\geq j\geq 1} \left(\frac{e}{T^i} \frac{e}{T^j}v_{n+1}\right)^* \bigg]
\nonumber \\
\hphantom{t^a}{}
\overset{\mu} \mapsto 	 \sum_{j=1}^n m_j z_j^a (v_1)^* \otimes \dots \otimes (f v_j)^* \otimes\dots\otimes (v_{n+1})^* \nonumber \\
\hphantom{t^a \overset{\eta^0}\mapsto}{} + (v_1)^* \otimes \dots \otimes (v_n)^* \otimes \bigg[\left(-m_{n+1}-2+a(k+2)\right) \left(\frac{e}{T^a}v_{n+1} \right)^* \nonumber \\
\hphantom{t^a \overset{\eta^0}\mapsto}{}+ \sum_{l=0}^{a-2} \bigg[{-} \frac{h}{T^{l+1}} \left(\frac{e}{T^{a-l-1}} v_{n+1} \right)^* + 2\frac{f}{T^l} \sum_{i+j=a-l\atop i\geq j\geq1} \left(\frac{e}{T^i} \frac{e}{T^j} v_{n+1} \right)^* \bigg] \bigg] \nonumber \\
\hphantom{t^a \overset{\eta^0}\mapsto}{}
-(v_1)^*\otimes\dots\otimes (v_n)^* \otimes \sum_{l=0}^{a-2} \bigg[ 2\frac{f}{T^l} \sum_{i+j=a-l\atop i\geq j\geq1} \left(\frac{e}{T^i} \frac{e}{T^j} v_{n+1} \right)^* - \frac{h}{T^{l+1}} \left(\frac{e}{T^{a-l-1}} v_{n+1} \right)^* \bigg] \nonumber \\
\hphantom{t^a \overset{\eta^0}\mapsto}{} - \sum_{s=1}^{a-1} \sum_{j=1}^n m_j z_j^s (v_1)^* \otimes \dots \otimes (v_n)^* \otimes \left(\frac{e}{T^{a-s}} v_{n+1}\right)^* \nonumber \\
\hphantom{t^a}{} = \bigg(a\kappa - \sum_{j=1}^n m_j\bigg) (v_1)^* \otimes \dots \otimes (v_n)^* \otimes \left(\frac{e}{T^a} v_{n+1}\right)^* \nonumber \\
\hphantom{t^a \overset{\eta^0}\mapsto}{} - \sum_{s=1}^{a-1} \sum_{j=1}^n m_j z_j^s (v_1)^* \otimes \dots \otimes (v_n)^* \otimes \left(\frac{e}{T^{a-s}} v_{n+1}\right)^* \nonumber \\
\hphantom{t^a \overset{\eta^0}\mapsto}{} + \sum_{j=1}^n m_j z_j^a (v_1)^* \otimes \dots \otimes (f v_j)^* \otimes\dots\otimes (v_{n+1})^*. \nonumber
\end{gather*}
In this calculation we use formula~(\ref{id B}) to express the action of $\frac{e}{T^a}$ on $(v_{n+1})^*$. Notice also that calculating the action on $V_{n+1}^*$ we use the automorphism~$\pi$, see Section~\ref{aut}.These formulas show that ${\rm d}\big(\eta^0(t^a)\big) = \eta^1(\partial(t^a))$.

Clearly the maps $\eta^1$, $\eta^2$ are injective. Theorem \ref{thm hom} is proved.
\end{proof}
